\documentclass[12pt, letterpaper]{article}
\usepackage{graphicx} % Required for inserting images
%\usepackage[T1]{fontenc}

\usepackage{amsfonts}
\usepackage{tabularx}
\usepackage{multirow}
\usepackage{booktabs}
\usepackage{float}
%\usepackage{natbib}

\usepackage{polski}
%\usepackage[utf8]{inputenc}

\title{Programy użytkowe \\ Lista 5}

\date{}

\begin{document}
\maketitle
Aby uruchomić Gnuplota w trybie interaktywnym, wprowadzamy w linii komend polecenie \texttt{gnuplot}.

Polecenia do tej listy wykonaj w trybie interaktywnym. Zapisz je w dokumencie \LaTeX{a}, dołączając do każdego polecenia wygenerowany przez nie wykres (w zadaniach, gdzie produkujemy kilka wykresów możesz użyć \verb|subfigure|). Wskazówka: do wyróżnienia kodu użyj polecenia \verb|\texttt{}|.
\begin{enumerate}
    \item Narysuj wykres funkcji $\sin{(x)},\, \sin{(2x)},\,\sin{(2x)}+\sin{(x)}$
    \begin{enumerate}
        \item każdy na osobnym rysunku,
        \item na wspólnym rysunku.
    \end{enumerate}
    Dodaj podpisy osi oraz tytuł i legendę.
    \item Wykonaj wykres funkcji $y = \tg{(x)}$ i $y = 1/x$. Zlokalizuj na nim dodatnie rozwiązanie równania $\tg{(x)} = 1/x$. Na podstawie wykresu znajdź przybliżoną wartość tego rozwiązania (z dokładnością do co najmniej 4 cyfr po przecinku). Wskazówka: w oknie terminala powiększ wykres w okolicach „podejrzanego” punktu. W lewym dolnym rogu wyświetlane są współrzędne miejsca, nad którym znajduje się kursor.
    \item Sprawdź jak wygląda wykres funkcji $\tg{(x)}$ w okolicy asymptoty. Wyjaśnij obserwowany efekt. Wskazówka: \texttt{set samples}.
    \item Nanieś na wykres dane z pliku \texttt{dane1.dat}. Użyj punktów, łamanej. Dopasuj krzywą gładką (wskazówka \verb|smooth csplines|).
\end{enumerate}
\end{document}